\documentclass[11pt]{article}

\usepackage[a4paper,margin=29mm]{geometry}
\usepackage[T1]{fontenc}
\usepackage[utf8]{inputenc}
\usepackage{lmodern}
\usepackage{microtype}
\usepackage{amsmath,amssymb,amsthm,mathtools}
\usepackage{booktabs,tabularx,array}
\usepackage{enumitem}
\usepackage{xcolor}
\usepackage[round,authoryear]{natbib}
\usepackage[colorlinks=true,linkcolor=blue!55!black,
  citecolor=red!45!black,urlcolor=blue!55!black]{hyperref}
\usepackage[nameinlink,noabbrev]{cleveref}

\hypersetup{
  pdftitle={The Biconvex Relaxation of Gromov--Wasserstein Transport},
  pdfauthor={Gabriel Peyre},
  pdfsubject={Measure-theoretic biconvex relaxations and alternating minimization for Gromov--Wasserstein transport},
  pdfkeywords={Gromov--Wasserstein, biconvex relaxation, alternating minimization, negative type, Sinkhorn}
}

\newtheorem{theorem}{Theorem}[section]
\newtheorem{proposition}[theorem]{Proposition}
\newtheorem{lemma}[theorem]{Lemma}
\newtheorem{corollary}[theorem]{Corollary}
\theoremstyle{definition}
\newtheorem{definition}[theorem]{Definition}
\newtheorem{example}[theorem]{Example}
\theoremstyle{remark}
\newtheorem{remark}[theorem]{Remark}

\newcommand{\R}{\mathbb R}
\newcommand{\one}{\mathbf 1}
\newcommand{\Mcal}{\mathcal M}
\newcommand{\Pcal}{\mathcal P}
\newcommand{\Kcal}{\mathcal K}
\newcommand{\Umat}{\mathrm U}
\newcommand{\Pmat}{\mathrm P}
\newcommand{\Qmat}{\mathrm Q}
\newcommand{\Cmat}{\mathrm C}
\newcommand{\Dmat}{\mathrm D}
\newcommand{\Dmatp}{\mathrm D'}
\newcommand{\Hmat}{\mathrm H}
\newcommand{\Jmat}{\mathrm J}
\newcommand{\Gmat}{\mathrm G}
\newcommand{\Gmatp}{\mathrm G'}
\newcommand{\Ee}{\mathcal E}
\newcommand{\Fcal}{\mathcal F}
\newcommand{\Bcal}{\mathcal B}
\newcommand{\Rcal}{\mathcal R}
\newcommand{\Ccal}{\mathcal C}
\newcommand{\Span}{\operatorname{span}}
\newcommand{\argmin}{\operatorname*{arg\,min}}
\newcommand{\dd}{\,\mathrm d}
\newcommand{\eps}{\varepsilon}
\newcommand{\KL}{\mathrm{KL}}
\newcommand{\Ent}{\mathrm{Ent}}
\newcommand{\ip}[2]{\left\langle #1,#2\right\rangle}
\newcommand{\norm}[1]{\left\lVert #1\right\rVert}

\title{\textbf{The Biconvex Relaxation of\\
Gromov--Wasserstein Transport}\\
\large A measure-theoretic account of tightness and alternating minimization}
\author{Gabriel Peyr\'e}
\date{July 2026}

\begin{document}
\maketitle

\begin{abstract}
A standard Gromov--Wasserstein algorithm freezes one occurrence of the
coupling in the quadratic distortion and optimizes the other.  This note
presents the underlying mechanism directly on a convex family of measures
\(\Kcal\subset\Pcal(\mathcal Z)\).  A symmetric bilinear form \(\Ee\) that
is nonpositive on feasible differences yields a tight two-measure relaxation:
replacing \(\Ee(\pi,\pi)\) by \(\Ee(\pi,\xi)\) changes neither the minimum
value nor its diagonal minimizers.  Squared Gromov--Wasserstein transport is
the running example, with \(\mathcal Z=\mathcal X\times\mathcal Y\) and
\(\Kcal=\Pi(\alpha,\beta)\); negative-type geometry supplies the required
sign.  The same framework includes both unregularized GW and its entropic
regularization.  Exact alternating minimization decreases the diagonal
objective.  For entropic GW on compact spaces, its decrease is exactly the
sum of a relative entropy and the nonnegative curvature defect of the GW
form.  The iterates remain in a uniformly positive, equicontinuous set of
coupling densities.  This entropy geometry gives sufficient decrease and a
projected-gradient error bound without introducing a coordinate Frobenius
norm or the much stronger total-variation topology.  An analytic Fredholm
{\L}ojasiewicz--Simon argument then proves finite length and convergence of
the complete measure-valued sequence to a stationary coupling.  An abstract
version of the argument applies to more general analytic biconvex energies.
\end{abstract}

\section{A quadratic problem over measures}
\label{sec:measure-setup}

The biconvex mechanism is most transparent before choosing coordinates or
discretizing the underlying spaces.  Let \(\mathcal Z\) be a compact metric
space, let \(\Mcal(\mathcal Z)\) be the vector space of finite signed Borel
measures, and let \(\Kcal\subset\Pcal(\mathcal Z)\) be a nonempty narrowly
compact convex set.  The two variables of the relaxation will be denoted
\(\pi,\xi\in\Kcal\).

Let
\(
  \Ee:\Mcal(\mathcal Z)\times\Mcal(\mathcal Z)\to\R
\)
be a symmetric bilinear form that is continuous on \(\Kcal\times\Kcal\).
The main example is generated by a bounded continuous symmetric kernel
\(L:\mathcal Z^2\to\R\):
\begin{equation}\label{eq:kernel-bilinear-form}
  \Ee(\pi,\xi)
  =
  \iint_{\mathcal Z^2}L(z,z')\dd\pi(z)\dd\xi(z').
\end{equation}
Freezing \(\xi\) produces the linear cost
\begin{equation}\label{eq:frozen-cost}
  \Ccal_\xi(z)
  :=\int_{\mathcal Z}L(z,z')\dd\xi(z'),
  \qquad
  \Ee(\pi,\xi)=\int_{\mathcal Z}\Ccal_\xi\dd\pi.
\end{equation}

\begin{definition}[Negative form on feasible differences]
\label{def:negative-form}
The bilinear form \(\Ee\) is \emph{negative on feasible differences} when
\begin{equation}\label{eq:negative-feasible-differences}
  \Ee(\sigma,\sigma)\leq0
  \qquad
  \text{for every }\sigma\in\Span(\Kcal-\Kcal).
\end{equation}
It is strictly negative when equality is possible only for \(\sigma=0\).
\end{definition}

This is conditional rather than global negativity: differences of
probability measures have zero total mass, and differences of transport plans
have zero marginals.  These linear conservation laws are precisely what makes
the GW form negative.

\paragraph{The running GW example.}
Let
\(
 \mathbb X=(\mathcal X,d_{\mathcal X},\alpha)
\)
and
\(
 \mathbb Y=(\mathcal Y,d_{\mathcal Y},\beta)
\)
be compact metric-measure spaces.  Set
\begin{equation}\label{eq:gw-measure-specialization}
  \mathcal Z=\mathcal X\times\mathcal Y,
  \qquad
  \Kcal=\Pi(\alpha,\beta),
  \qquad
  L((x,y),(x',y'))
  =\bigl(d_{\mathcal X}(x,x')-d_{\mathcal Y}(y,y')\bigr)^2.
\end{equation}
Then
\begin{equation}\label{eq:continuous-gw-energy}
  \operatorname{GW}_2^2(\mathbb X,\mathbb Y)
  =\min_{\pi\in\Pi(\alpha,\beta)}\Ee(\pi,\pi).
\end{equation}
For discrete measures
\(
 \alpha=\sum_i a_i\delta_{x_i}
\)
and
\(
 \beta=\sum_j b_j\delta_{y_j},
\)
a coupling is represented by
\(
 \Pmat\in\Umat(\mathrm a,\mathrm b)
\)
and \eqref{eq:kernel-bilinear-form} becomes
\begin{equation}\label{eq:discrete-polarized-gw}
  \Ee(\Pmat,\Qmat)
  =
  \sum_{i,i',j,j'}
  \bigl(D_{ii'}-D'_{jj'}\bigr)^2P_{ij}Q_{i'j'}.
\end{equation}
Here \(\mathrm a=(a_i)_i\), \(\mathrm b=(b_j)_j\),
\(D_{ii'}=d_{\mathcal X}(x_i,x_{i'})\), and
\(D'_{jj'}=d_{\mathcal Y}(y_j,y_{j'})\).  The measure notation and matrix
notation thus describe the same object.

\section{Tightness of the biconvex relaxation}
\label{sec:tightness}

Duplicating the measure always lowers the minimum because diagonal pairs
remain feasible.  Conditional negativity is exactly what prevents the lower
bound from becoming strict.  The result is the measure-level form of the
quadratic principle of \citet{Konno1976}, used for generalized GW by
\citet{SejourneVialardPeyre2021}.

\begin{theorem}[Tight biconvex relaxation over measures]
\label{thm:measure-tightness}
Let \(\Rcal:\Kcal\to\R\cup\{+\infty\}\) be proper and narrowly lower
semicontinuous, and define
\begin{align}
  \Fcal(\pi)
  &=\Ee(\pi,\pi)+\Rcal(\pi),
  \label{eq:general-diagonal-objective}\\
  \Bcal(\pi,\xi)
  &=\Ee(\pi,\xi)
    +\frac12\Rcal(\pi)+\frac12\Rcal(\xi).
  \label{eq:general-biconvex-objective}
\end{align}
Assume \eqref{eq:negative-feasible-differences}.  Then both minima are
attained and
\begin{equation}\label{eq:measure-tight-values}
  \min_{\pi,\xi\in\Kcal}\Bcal(\pi,\xi)
  =
  \min_{\pi\in\Kcal}\Fcal(\pi).
\end{equation}
If \((\pi_*,\xi_*)\) minimizes \(\Bcal\), then
\begin{equation}\label{eq:measure-minimizer-characterization}
  \Fcal(\pi_*)=\Fcal(\xi_*)=\min_{\Kcal}\Fcal,
  \qquad
  \Ee(\pi_*-\xi_*,\pi_*-\xi_*)=0.
\end{equation}
Every relaxed minimizer is diagonal if \(\Ee\) is strictly negative on
nonzero feasible differences or if \(\Rcal\) is strictly convex on its
finite domain.  If \(\Rcal\) is convex, \(\Bcal\) is convex in each measure
separately.
\end{theorem}

\begin{proof}
Compactness and lower semicontinuity give existence.  For every
\(\pi,\xi\in\Kcal\), polarization gives
\begin{equation}\label{eq:measure-polarization-identity}
  \Bcal(\pi,\xi)
  =
  \frac12\Fcal(\pi)+\frac12\Fcal(\xi)
  -\frac12\Ee(\pi-\xi,\pi-\xi).
\end{equation}
The sign condition implies that the right-hand side is at least
\(
 [\Fcal(\pi)+\Fcal(\xi)]/2
\), hence at least \(\min_{\Kcal}\Fcal\).  Conversely,
\(\Bcal(\pi,\pi)=\Fcal(\pi)\), which proves
\eqref{eq:measure-tight-values}.  Equality in the two inequalities gives
\eqref{eq:measure-minimizer-characterization}.

Strict negativity immediately forces \(\pi_*=\xi_*\).  If instead
\(\Rcal\) is strictly convex and \(\pi_*\ne\xi_*\), then
\eqref{eq:measure-minimizer-characterization} makes the quadratic part affine
on the segment joining them, while \(\Rcal\) is strictly below its chord.
The midpoint would have a smaller value of \(\Fcal\), a contradiction.
Separate convexity follows because the bilinear term is linear in either
measure.
\end{proof}

\begin{example}[Why negativity is necessary]
Take \(\mathcal Z=\{-1,1\}\), identify a signed measure with its first
moment \(x\in[-1,1]\), and set \(\Ee(x,y)=xy\), \(\Rcal=0\).  The diagonal
minimum of \(x^2\) is zero, whereas the two-variable minimum of \(xy\) is
\(-1\).  Duplicating a variable is not a tight relaxation for a positive
quadratic form.
\end{example}

\paragraph{Concavity and tangent majorization.}
The same sign condition has an algorithmic consequence.  The diagonal
quadratic map is concave along feasible segments and, for all
\(\pi,\xi\in\Kcal\),
\begin{equation}\label{eq:measure-tangent-majorization}
  \Ee(\xi,\xi)
  \leq
  2\Ee(\pi,\xi)-\Ee(\pi,\pi).
\end{equation}
Indeed, this is exactly the inequality
\(
 \Ee(\xi-\pi,\xi-\pi)\leq0
\)
written after expansion.  Thus freezing one occurrence of the measure gives
a tangent upper bound, not merely one block of a biconvex objective.

\section{Why squared GW has the required sign}
\label{sec:negative-type}

The abstract condition becomes geometric for squared GW.  Recall that a
symmetric kernel \(d\) on \(\mathcal X\) is of negative type when
\begin{equation}\label{eq:negative-type-kernel}
  \iint d(x,x')\dd\eta(x)\dd\eta(x')\leq0
  \qquad
  \text{whenever }\eta(\mathcal X)=0.
\end{equation}
By Schoenberg's theorem, this is equivalent, after the usual harmless
normalization on the diagonal, to the existence of a Hilbert space
\(\mathcal H\) and a map \(\varphi:\mathcal X\to\mathcal H\) such that
\(
 d(x,x')=\norm{\varphi(x)-\varphi(x')}^2
\)
\citep{Schoenberg1938}.  Euclidean powers
\(d(x,x')=\norm{x-x'}^r\), \(0<r\leq2\), and tree distances are standard
examples.

\begin{proposition}[Conditional negativity of the squared GW form]
\label{prop:measure-gw-negative}
Assume that \(d_{\mathcal X}\) and \(d_{\mathcal Y}\) are kernels of
negative type.  For every signed measure \(\sigma\) on
\(\mathcal X\times\mathcal Y\) with zero marginals,
\begin{equation}\label{eq:measure-gw-negative-square}
  \Ee(\sigma,\sigma)
  =
  -8
  \left\|
    \int
    \varphi(x)\otimes\psi(y)\dd\sigma(x,y)
  \right\|_{\mathcal H_{\mathcal X}\otimes\mathcal H_{\mathcal Y}}^2
  \leq0,
\end{equation}
where \(\varphi\) and \(\psi\) are Hilbert embeddings of the two kernels.
Consequently, the GW specialization \eqref{eq:gw-measure-specialization}
satisfies Definition~\ref{def:negative-form}.
\end{proposition}

\begin{proof}
Expand the squared difference in \eqref{eq:gw-measure-specialization}.  The
terms involving only \(d_{\mathcal X}^2\) or only
\(d_{\mathcal Y}^2\) vanish because \(\sigma\) has zero marginals.  Hence
\begin{equation}\label{eq:measure-gw-cross-term}
  \Ee(\sigma,\sigma)
  =-2\iint
  d_{\mathcal X}(x,x')d_{\mathcal Y}(y,y')
  \dd\sigma(x,y)\dd\sigma(x',y').
\end{equation}
Insert the two squared-Hilbert-distance representations.  Every term
containing a function of \(x\) alone or \(y\) alone again vanishes by the
zero-marginal constraints.  The only surviving term is four times the squared
Hilbert--Schmidt norm in \eqref{eq:measure-gw-negative-square}; multiplication
by \(-2\) proves the claim.  Differences of couplings in
\(\Pi(\alpha,\beta)\) have zero marginals, so the last statement follows.
\end{proof}

\paragraph{Finite-dimensional coordinates.}
For finite spaces, let \(\Dmat,\Dmatp\) be the two distance matrices and let
\(
 \Jmat_n=I_n-\one_n\one_n^\top/n
\), with \(\Jmat_m\) defined similarly.  The centered Gram matrices
\begin{equation}\label{eq:discrete-centered-gram}
  \Gmat=-\frac12\Jmat_n\Dmat\Jmat_n,
  \qquad
  \Gmatp=-\frac12\Jmat_m\Dmatp\Jmat_m
\end{equation}
are positive semidefinite.  A zero-marginal signed matrix \(\Hmat\) satisfies
\begin{equation}\label{eq:discrete-gw-negative-square}
  \Ee(\Hmat,\Hmat)
  =-8\norm{\Gmat^{1/2}\Hmat\Gmatp^{1/2}}_{\mathrm F}^2.
\end{equation}
This is the coordinate version of
\eqref{eq:measure-gw-negative-square}.  It also shows when strictness fails:
different couplings can be invisible to one of the centered geometries.

\begin{corollary}[Tight squared-GW relaxation]
\label{cor:gw-tightness}
Under the hypotheses of Proposition~\ref{prop:measure-gw-negative},
\begin{equation}\label{eq:continuous-gw-tightness}
  \min_{\pi,\xi\in\Pi(\alpha,\beta)}\Ee(\pi,\xi)
  =
  \min_{\pi\in\Pi(\alpha,\beta)}\Ee(\pi,\pi).
\end{equation}
Each component of a relaxed minimizer is a GW minimizer.  If the tensor
feature map in \eqref{eq:measure-gw-negative-square} is injective on
zero-marginal signed measures, every relaxed minimizer is diagonal.
\end{corollary}

The conclusion is specific to the squared outer loss and to compatible signs
of the two within-space kernels.  For other distortions the polarized form
remains linear in each coupling but may fail to be negative on feasible
differences; the resulting two-plan problem is then only a lower bound.

\section{Entropic and unregularized GW in one formulation}
\label{sec:entropic-gw}

Choose a reference probability \(\mathfrak m\in\Pcal(\mathcal Z)\) and use
the lower-semicontinuous convention
\begin{equation}\label{eq:relative-entropy}
  \KL(\pi\mid\mathfrak m)
  =
  \begin{cases}
  \displaystyle
  \int\bigl(\rho\log\rho-\rho+1\bigr)\dd\mathfrak m,
  &\pi=\rho\mathfrak m,\\
  +\infty,&\text{otherwise}.
  \end{cases}
\end{equation}
For \(\eps\geq0\), set
\begin{equation}\label{eq:epsilon-regularizer}
  \Rcal_\eps(\pi)
  =
  \begin{cases}
    0,&\eps=0,\\
    \eps\KL(\pi\mid\mathfrak m),&\eps>0,
  \end{cases}
\end{equation}
so that no ambiguous product \(0\cdot(+\infty)\) occurs.  Define
\begin{align}
  \Fcal_\eps(\pi)
  &=\Ee(\pi,\pi)+\Rcal_\eps(\pi),
  \label{eq:epsilon-diagonal-objective}\\
  \Bcal_\eps(\pi,\xi)
  &=\Ee(\pi,\xi)
  +\frac12\Rcal_\eps(\pi)
  +\frac12\Rcal_\eps(\xi).
  \label{eq:epsilon-biconvex-objective}
\end{align}
The factors \(1/2\) ensure
\(
 \Bcal_\eps(\pi,\pi)=\Fcal_\eps(\pi)
\).

\begin{corollary}[Tightness with and without entropy]
\label{cor:entropic-tightness}
Under \eqref{eq:negative-feasible-differences},
\begin{equation}\label{eq:epsilon-tightness}
  \min_{\pi,\xi\in\Kcal}\Bcal_\eps(\pi,\xi)
  =
  \min_{\pi\in\Kcal}\Fcal_\eps(\pi)
  \qquad(\eps\geq0).
\end{equation}
For \(\eps>0\), every finite-valued relaxed minimizer is diagonal.  For
\(\eps=0\), its two components are diagonal minimizers but need not coincide.
\end{corollary}

For GW, the canonical choice is
\begin{equation}\label{eq:gw-reference-product}
  \mathfrak m=\alpha\otimes\beta,
  \qquad
  \Kcal=\Pi(\alpha,\beta).
\end{equation}
Then \(\eps=0\) gives ordinary GW and \(\eps>0\) gives entropic GW.  In the
discrete setting, \(\KL(\pi\mid\alpha\otimes\beta)\) differs from the usual
negative-entropy penalty on the coupling matrix only by a constant determined
by the prescribed marginals.

\section{Alternating minimization}
\label{sec:alternating}

The lifted problem suggests a one-sequence algorithm.  Starting from
\(\pi^{(0)}\in\Kcal\), define
\begin{equation}\label{eq:measure-alternating-update}
  \pi^{(\ell+1)}
  \in
  \argmin_{\pi\in\Kcal}
  \left\{
    \Ee(\pi^{(\ell)},\pi)
    +\frac12\Rcal_\eps(\pi)
  \right\}.
\end{equation}
This is exact alternating minimization of \(\Bcal_\eps\), with the names of
the two symmetric blocks exchanged after each update.  It is also a
majorization--minimization step for \(\Fcal_\eps\).

\paragraph{Specialization to GW.}
For \(\Kcal=\Pi(\alpha,\beta)\), the frozen cost is
\begin{equation}\label{eq:gw-frozen-measure-cost}
  \Ccal_{\pi^{(\ell)}}(x,y)
  =
  \int
  \bigl(d_{\mathcal X}(x,x')-d_{\mathcal Y}(y,y')\bigr)^2
  \dd\pi^{(\ell)}(x',y').
\end{equation}
Thus \eqref{eq:measure-alternating-update} is an ordinary OT problem with
this cost.  It is a linear Kantorovich problem when \(\eps=0\) and an
entropic OT problem, solvable by Sinkhorn, when \(\eps>0\).  The regularization
strength of this block problem is \(\eps/2\), not \(\eps\), because each of
the two lifted variables carries half of the diagonal entropy.  In matrix
coordinates the cost is
\begin{equation}\label{eq:discrete-frozen-gw-cost}
  \Cmat(\Pmat)
  =
  \Dmat^{\odot2}\mathrm a\one_m^\top
  +\one_n(\Dmatp^{\odot2}\mathrm b)^\top
  -2\Dmat\Pmat\Dmatp.
\end{equation}

\begin{proposition}[Descent of the diagonal objective]
\label{prop:measure-descent}
Assume \eqref{eq:negative-feasible-differences}.  Every exact update
\eqref{eq:measure-alternating-update} satisfies
\begin{equation}\label{eq:measure-objective-descent}
  \Fcal_\eps(\pi^{(\ell+1)})
  \leq
  \Fcal_\eps(\pi^{(\ell)}).
\end{equation}
This holds for every \(\eps\geq0\), with any selected block minimizer when
\(\eps=0\).
\end{proposition}

\begin{proof}
Put \(\pi=\pi^{(\ell)}\) and \(\xi=\pi^{(\ell+1)}\).  The tangent bound
\eqref{eq:measure-tangent-majorization} gives
\(
 \Ee(\xi,\xi)\leq2\Ee(\pi,\xi)-\Ee(\pi,\pi)
\).
Optimality of the block update, after multiplying by two, gives
\(
 2\Ee(\pi,\xi)+\Rcal_\eps(\xi)
 \leq
 2\Ee(\pi,\pi)+\Rcal_\eps(\pi)
\).
Combining the inequalities proves \eqref{eq:measure-objective-descent}.
\end{proof}

\paragraph{Fixed points and stationarity.}
Let \(N_{\Kcal}(\pi)\) denote the convex normal cone and interpret
subdifferentials in the ambient signed-measure space.  If \(\pi_*\) is a
fixed point of \eqref{eq:measure-alternating-update}, its block optimality
condition, multiplied by two, is
\begin{equation}\label{eq:measure-stationarity}
  0
  \in
  2\Ccal_{\pi_*}
  +\partial\Rcal_\eps(\pi_*)
  +N_{\Kcal}(\pi_*).
\end{equation}
This is exactly first-order stationarity for \(\Fcal_\eps\).  In the
entropic GW case, writing
\(
 \rho_*=\dd\pi_*/\dd(\alpha\otimes\beta)
\), it takes the variational form
\begin{equation}\label{eq:entropic-gw-variational-stationarity}
  \int
  \bigl(2\Ccal_{\pi_*}+\eps\log\rho_*\bigr)
  \dd(\pi-\pi_*)
  \geq0
  \qquad
  \text{for every }\pi\in\Pi(\alpha,\beta).
\end{equation}
Equivalently, under the usual positivity assumptions, the expression in
parentheses equals a sum of one potential on \(\mathcal X\) and one potential
on \(\mathcal Y\), almost everywhere.

\section{Intrinsic entropy descent on coupling measures}
\label{sec:convergence}

The useful geometry is already present in the regularizer.  We therefore
measure one block step by its entropy Bregman divergence rather than by a
coordinate norm.  Specialize here to
\(\Kcal=\Pi(\alpha,\beta)\), put
\(\mathfrak m=\alpha\otimes\beta\), and identify
\(\pi=\rho\mathfrak m\) with its density.  The feasible densities and their
linear tangent space are
\begin{align}
 \mathcal A
 &=
 \left\{
 \rho\geq0:
 \int\rho(x,y)\dd\beta(y)=1,\ 
 \int\rho(x,y)\dd\alpha(x)=1
 \right\},
 \label{eq:feasible-coupling-densities}\\
 \mathsf T
 &=
 \left\{
 h\in L^2(\mathfrak m):
 \int h(x,y)\dd\beta(y)=0,\ 
 \int h(x,y)\dd\alpha(x)=0
 \right\}.
 \label{eq:coupling-density-tangent}
\end{align}
All equalities are understood almost everywhere.  Let \(K\) be the integral
operator induced by \(L\),
\begin{equation}\label{eq:gw-integral-operator}
 (K\rho)(z)=\int_{\mathcal Z}L(z,z')\rho(z')\dd\mathfrak m(z'),
 \qquad
 \Ee(\rho\mathfrak m,\eta\mathfrak m)
 =\ip{\rho}{K\eta}_{L^2(\mathfrak m)}.
\end{equation}
Writing
\(\Ent(\rho)=\int(\rho\log\rho-\rho+1)\dd\mathfrak m\), the diagonal
functional becomes
\begin{equation}\label{eq:entropic-gw-density-functional}
 \Phi_\eps(\rho)
 =
 \ip{\rho}{K\rho}_{L^2(\mathfrak m)}
 +\eps\Ent(\rho).
\end{equation}

\subsection{An exact Bregman descent identity}

The sign of the polarized GW form does more than prove monotonicity: it
separates the decrease into two nonnegative terms.

\begin{proposition}[Exact entropy descent]
\label{prop:exact-entropy-descent}
Let \(\eps>0\), let \(\rho^{(\ell)},\rho^{(\ell+1)}\in\mathcal A\) be
positive densities, and suppose that \(\rho^{(\ell+1)}\) is the block update
\eqref{eq:measure-alternating-update}.  Then
\begin{align}
 \Phi_\eps(\rho^{(\ell)})-\Phi_\eps(\rho^{(\ell+1)})
 &=
 \eps\KL\!\left(
   \rho^{(\ell)}\mathfrak m
   \,\middle|\,
   \rho^{(\ell+1)}\mathfrak m
 \right)
 \notag\\
 &\quad
 -\Ee\!\left(
   (\rho^{(\ell+1)}-\rho^{(\ell)})\mathfrak m,
   (\rho^{(\ell+1)}-\rho^{(\ell)})\mathfrak m
 \right).
 \label{eq:exact-entropic-gw-descent}
\end{align}
Both terms on the right are nonnegative.
\end{proposition}

\begin{proof}
Put \(\rho=\rho^{(\ell)}\), \(\widehat\rho=\rho^{(\ell+1)}\), and
\(d=\widehat\rho-\rho\in\mathsf T\).  Since the constraints in
\eqref{eq:feasible-coupling-densities} are affine and the entropic block
minimizer is positive, its first-order condition in the feasible direction
\(\rho-\widehat\rho\) is an equality:
\[
 \ip{K\rho+\frac\eps2\log\widehat\rho}
     {\rho-\widehat\rho}_{L^2(\mathfrak m)}
 =0.
\]
The Bregman identity for negative entropy therefore gives
\[
 \ip{\rho}{K\rho}+\frac\eps2\Ent(\rho)
 -\ip{\rho}{K\widehat\rho}-\frac\eps2\Ent(\widehat\rho)
 =
 \frac\eps2\KL(\rho\mathfrak m\mid\widehat\rho\mathfrak m).
\]
The polarization identity now reads
\[
 \Phi_\eps(\rho)-\Phi_\eps(\widehat\rho)
 =
 2\left[
  \ip{\rho}{K\rho}+\frac\eps2\Ent(\rho)
  -\ip{\rho}{K\widehat\rho}-\frac\eps2\Ent(\widehat\rho)
 \right]
 -\ip{d}{Kd}.
\]
Substituting the preceding equality proves
\eqref{eq:exact-entropic-gw-descent}.  Finally,
\(-\ip{d}{Kd}=-\Ee(d\mathfrak m,d\mathfrak m)\geq0\) because
\(d\mathfrak m\) is a feasible difference.
\end{proof}

Since \(L\) is bounded and entropy is nonnegative, the objective is bounded
below.  Thus its values converge and
\begin{equation}\label{eq:summable-intrinsic-descent}
 \sum_{\ell\geq0}
 \KL\!\left(
   \rho^{(\ell)}\mathfrak m
   \,\middle|\,
   \rho^{(\ell+1)}\mathfrak m
 \right)
 <+\infty.
\end{equation}
This conclusion is intrinsic to entropy.  No total-variation estimate is
needed.

\subsection{A compact positive invariant set}

For a bounded cost kernel, every entropic block problem is a Schr\"odinger
scaling problem with a uniformly positive Gibbs kernel.  This prevents the
iterates from approaching the boundary of the positive cone.

\begin{proposition}[Uniform positivity and compactness of GW updates]
\label{prop:uniform-positive-gw-updates}
Assume that \(\mathcal X,\mathcal Y\) are compact, that
\(\alpha,\beta\) have full support, and that \(L\) is continuous.  Set
\(M=\norm{L}_\infty\) and
\[
 k_-=\exp(-2M/\eps),
 \qquad
 k_+=\exp(2M/\eps),
 \qquad
 c_\eps=(k_-/k_+)^2,
 \qquad
 C_\eps=(k_+/k_-)^2.
\]
Every block update has a continuous density satisfying
\begin{equation}\label{eq:uniform-density-box}
 0<c_\eps\leq\rho^{(\ell)}(x,y)\leq C_\eps<+\infty
 \qquad(\ell\geq1).
\end{equation}
The family \((\rho^{(\ell)})_{\ell\geq1}\) is equicontinuous and hence
relatively compact in \(C(\mathcal X\times\mathcal Y)\).
\end{proposition}

\begin{proof}
Since \(\pi^{(\ell)}\) is a probability measure,
\(\norm{\Ccal_{\pi^{(\ell)}}}_\infty\leq M\).  The block Gibbs kernel
\[
 k_\ell(x,y)
 =
 \exp\!\left(-\frac{2\Ccal_{\pi^{(\ell)}}(x,y)}{\eps}\right)
\]
therefore lies between \(k_-\) and \(k_+\).  The unique block minimizer has
the scaling form
\(
 \rho^{(\ell+1)}(x,y)=u_\ell(x)k_\ell(x,y)v_\ell(y)
\).
Use the scalar indeterminacy of the two scalings to impose
\(\int v_\ell\dd\beta=1\).  The two marginal equations then give
\[
 k_+^{-1}\leq u_\ell\leq k_-^{-1},
 \qquad
 k_-/k_+\leq v_\ell\leq k_+/k_-,
\]
which proves \eqref{eq:uniform-density-box}.

Uniform continuity of \(L\) gives a common modulus of continuity for all
frozen costs, hence for all \(k_\ell\).  The formulas
\[
 u_\ell(x)=
 \left(\int k_\ell(x,y)v_\ell(y)\dd\beta(y)\right)^{-1},
 \qquad
 v_\ell(y)=
 \left(\int k_\ell(x,y)u_\ell(x)\dd\alpha(x)\right)^{-1}
\]
and the preceding lower bounds transfer this modulus to the scalings and to
\(\rho^{(\ell+1)}\).  Arzel\`a--Ascoli gives relative compactness.
\end{proof}

On the box \eqref{eq:uniform-density-box}, Taylor's formula for
\(s\mapsto s\log s-s+1\) yields
\begin{equation}\label{eq:kl-fisher-equivalence}
 \frac1{2C_\eps}\norm{\rho-\eta}_{L^2(\mathfrak m)}^2
 \leq
 \KL(\rho\mathfrak m\mid\eta\mathfrak m)
 \leq
 \frac1{2c_\eps}\norm{\rho-\eta}_{L^2(\mathfrak m)}^2.
\end{equation}
Consequently, \eqref{eq:summable-intrinsic-descent} implies
square-summability of the density increments in \(L^2(\mathfrak m)\).

\begin{corollary}[Stationarity of cluster points]
\label{cor:stationary-gw-cluster-points}
Under the assumptions of
Proposition~\ref{prop:uniform-positive-gw-updates}, the objective values
converge, the density increments tend to zero in \(L^2(\mathfrak m)\), and
every uniform cluster point is a fixed point of the block map and satisfies
\eqref{eq:entropic-gw-variational-stationarity}.
\end{corollary}

\begin{proof}
Only the last claim remains to prove.  Let
\(\rho^{(\ell_k)}\to\rho_*\) uniformly.  Relative compactness gives a
uniformly convergent subsequence of \(\rho^{(\ell_k+1)}\); its limit equals
\(\rho_*\), because the corresponding \(L^2\) increments tend to zero and
\(\mathfrak m\) has full support.  For every \(\eta\in\mathcal A\) of finite
entropy, block optimality gives
\[
 \ip{\rho^{(\ell_k)}}{K\rho^{(\ell_k+1)}}
 +\frac\eps2\Ent(\rho^{(\ell_k+1)})
 \leq
 \ip{\rho^{(\ell_k)}}{K\eta}
 +\frac\eps2\Ent(\eta).
\]
For infinite-entropy competitors the limiting inequality is automatic.  The
uniform density bounds make entropy continuous along this subsequence,
while boundedness of \(K\) gives continuity of the bilinear terms.  Passing
to the limit shows that \(\rho_*\) uniquely minimizes the block functional
frozen at \(\rho_*\).  Thus it is a fixed point.  Since it is uniformly
positive, its affine first-order condition is
\[
 \int
 \bigl(2K\rho_*+\eps\log\rho_*\bigr)h\dd\mathfrak m=0
 \qquad(h\in\mathsf T),
\]
which is equivalent to
\eqref{eq:entropic-gw-variational-stationarity}.
\end{proof}

This still does not by itself prove convergence of the full sequence:
square-summable increments need not have summable lengths.  The next section
supplies the missing finite-length mechanism.

\section{Full convergence of alternating minimization}
\label{sec:full-alternating-convergence}

The final step is a measure-space analogue of the
Kurdyka--{\L}ojasiewicz argument used for finite-dimensional block
minimization \citep{AttouchBolteRedontSoubeyran2010,XuYin2013}.  The correct
local quadratic structure is the Fisher metric generated by entropy, not a
Frobenius norm attached to a discretization.  At a positive density \(\rho\),
it is
\begin{equation}\label{eq:fisher-tangent-norm}
 \norm{h}_{\rho,\mathrm{F}}^2
 =
 \int\frac{h^2}{\rho}\dd\mathfrak m,
 \qquad h\in\mathsf T.
\end{equation}
The bounds \eqref{eq:uniform-density-box} make this norm uniformly equivalent
to \(L^2(\mathfrak m)\) along the iterates.  The proof has three ingredients:
the exact entropy decrease controls the squared step, block optimality
controls the residual gradient by the preceding step, and an
infinite-dimensional {\L}ojasiewicz--Simon inequality turns these two
quadratic estimates into summable lengths.

\subsection{Projected gradient and relative error}

Let \(\mathsf P_{\mathsf T}\) be the \(L^2(\mathfrak m)\)-orthogonal
projection onto \(\mathsf T\).  It has the explicit bounded extension to
\(L^\infty(\mathfrak m)\)
\begin{align}
 (\mathsf P_{\mathsf T}q)(x,y)
 &=
 q(x,y)
 -\int q(x,y')\dd\beta(y')
 -\int q(x',y)\dd\alpha(x')
 \notag\\
 &\quad
 +\iint q(x',y')\dd\alpha(x')\dd\beta(y').
 \label{eq:coupling-tangent-projection}
\end{align}
The constrained gradient of \(\Phi_\eps\) is represented by
\begin{equation}\label{eq:projected-entropic-gw-gradient}
 G_\eps(\rho)
 =
 \mathsf P_{\mathsf T}\bigl(2K\rho+\eps\log\rho\bigr).
\end{equation}

\begin{proposition}[Block residual estimate]
\label{prop:block-residual-estimate}
The entropic GW update satisfies
\begin{align}
 \mathsf P_{\mathsf T}
 \left(K\rho^{(\ell)}
       +\frac\eps2\log\rho^{(\ell+1)}\right)
 &=0,
 \label{eq:block-projected-optimality}\\
 G_\eps(\rho^{(\ell+1)})
 &=
 2\mathsf P_{\mathsf T}K
 \bigl(\rho^{(\ell+1)}-\rho^{(\ell)}\bigr).
 \label{eq:block-relative-error-identity}
\end{align}
In particular, for a constant \(b\) depending only on \(L\),
\begin{equation}\label{eq:block-relative-error-bound}
 \norm{G_\eps(\rho^{(\ell+1)})}_{L^\infty(\mathfrak m)}
 \leq
 b\norm{\rho^{(\ell+1)}-\rho^{(\ell)}}_{L^2(\mathfrak m)}.
\end{equation}
\end{proposition}

\begin{proof}
Equation \eqref{eq:block-projected-optimality} is the first-order condition
of the affine-constrained block problem.  Subtracting twice this identity
from \eqref{eq:projected-entropic-gw-gradient} at
\(\rho^{(\ell+1)}\) gives
\eqref{eq:block-relative-error-identity}.  Finally,
\(\mathsf P_{\mathsf T}\) is bounded on \(L^\infty\), while
\[
 \norm{Kh}_\infty
 \leq\norm{L}_\infty\norm{h}_{L^1(\mathfrak m)}
 \leq\norm{L}_\infty\norm{h}_{L^2(\mathfrak m)}.
\]
\end{proof}

\subsection{The finite-length principle}

The following theorem isolates the convergence mechanism.  Its two ambient
spaces play different roles: compactness is required in a strong Banach
topology, whereas the length is measured in a weaker Hilbert norm.

\begin{theorem}[Measure-space {\L}ojasiewicz--Simon convergence principle]
\label{thm:measure-ls-convergence-principle}
Let \(\mathsf B\) be a Banach space continuously and injectively embedded
in a Hilbert space \(\mathsf H\), and let \(\mathsf G\) be a normed space.
Let \((\rho^{(\ell)})_\ell\) be relatively compact in \(\mathsf B\), let
\(\Phi\) be continuous on the \(\mathsf B\)-closure of the iterates, and
let \(\Phi(\rho^{(\ell)})\) decrease to \(\Phi_*\).  Assume that a residual
map \(G\) and constants \(a,b>0\) satisfy
\begin{align}
 \Phi(\rho^{(\ell)})-\Phi(\rho^{(\ell+1)})
 &\geq
 a\norm{\rho^{(\ell+1)}-\rho^{(\ell)}}_{\mathsf H}^2,
 \label{eq:abstract-sufficient-decrease}\\
 \norm{G(\rho^{(\ell+1)})}_{\mathsf G}
 &\leq
 b\norm{\rho^{(\ell+1)}-\rho^{(\ell)}}_{\mathsf H}.
 \label{eq:abstract-relative-error}
\end{align}
Suppose that every cluster point \(\rho_*\) is stationary and satisfies a
{\L}ojasiewicz--Simon inequality: there are a \(\mathsf B\)-neighborhood
\(U_*\), a constant \(c_*>0\), and \(\theta_*\in(0,1/2]\) such that
\begin{equation}\label{eq:lojasiewicz-simon-density}
 \left|\Phi(\rho)-\Phi(\rho_*)\right|^{1-\theta_*}
 \leq
 c_*\norm{G(\rho)}_{\mathsf G}
 \qquad(\rho\in U_*).
\end{equation}
Then the sequence has finite length,
\begin{equation}\label{eq:measure-finite-length}
 \sum_{\ell=0}^{+\infty}
 \norm{\rho^{(\ell+1)}-\rho^{(\ell)}}_{\mathsf H}<+\infty,
\end{equation}
and converges in \(\mathsf H\) to one stationary point.
\end{theorem}

\begin{proof}
Choose a cluster point \(\rho_*\), with convergence along a subsequence in
\(\mathsf B\).  Continuity of \(\Phi\) gives
\(\Phi(\rho_*)=\Phi_*\).  On the compact \(\mathsf B\)-closure
\(\mathfrak K\) of the iterates, the embedding into \(\mathsf H\) is a
continuous injection into a Hausdorff space.  Its inverse on its image is
therefore continuous.  Consequently, there is \(\delta>0\) such that
\begin{equation}\label{eq:strong-neighborhood-from-hilbert-ball}
 \rho\in\mathfrak K,
 \quad
 \norm{\rho-\rho_*}_{\mathsf H}<\delta
 \quad\Longrightarrow\quad
 \rho\in U_*.
\end{equation}
Write
\[
 e_\ell=\Phi(\rho^{(\ell)})-\Phi_*,
 \qquad
 d_\ell=\norm{\rho^{(\ell+1)}-\rho^{(\ell)}}_{\mathsf H},
 \qquad
 \varphi(s)=\frac{c_*}{\theta_*}s^{\theta_*}.
\]
Whenever \(\rho^{(\ell)}\in U_*\), \(e_\ell>0\), and
\(d_{\ell-1}>0\), concavity of \(\varphi\),
\eqref{eq:abstract-sufficient-decrease},
\eqref{eq:lojasiewicz-simon-density} at \(\rho^{(\ell)}\), and
\eqref{eq:abstract-relative-error} at the preceding step give
\[
 \varphi(e_\ell)-\varphi(e_{\ell+1})
 \geq
 \varphi'(e_\ell)(e_\ell-e_{\ell+1})
 \geq
 \frac{a\,d_\ell^2}{b\,d_{\ell-1}}.
\]
If \(d_{\ell-1}=0\), the relative-error and
{\L}ojasiewicz--Simon inequalities imply \(e_\ell=0\).  If
\(e_\ell=0\), monotonicity and sufficient decrease give
\(e_{\ell+1}=d_\ell=0\), and induction makes the remaining tail stationary.
Otherwise, the arithmetic--geometric mean inequality gives
\[
 2d_\ell
 \leq
 d_{\ell-1}
 +
 \frac ba
 \bigl(\varphi(e_\ell)-\varphi(e_{\ell+1})\bigr).
\]
Because \(d_\ell\to0\) by sufficient decrease and \(e_\ell\to0\), choose
an index \(n\) on the cluster subsequence such that
\begin{equation}\label{eq:ls-trapping-choice}
 \norm{\rho^{(n)}-\rho_*}_{\mathsf H}
 +d_{n-1}+\frac ba\varphi(e_n)
 <\delta.
\end{equation}
Suppose, by contradiction, that the tail first exits \(U_*\) at
\(\rho^{(J+1)}\).  The preceding estimate is valid for
\(\ell=n,\ldots,J\).  Summing it and cancelling common increments gives
\begin{equation}\label{eq:ls-finite-tail-bound}
 \sum_{\ell=n}^{J}d_\ell+d_J
 \leq
 d_{n-1}+\frac ba
 \bigl(\varphi(e_n)-\varphi(e_{J+1})\bigr)
 \leq
 d_{n-1}+\frac ba\varphi(e_n).
\end{equation}
Together with \eqref{eq:ls-trapping-choice}, this places
\(\rho^{(J+1)}\) within \(\delta\) of \(\rho_*\) in \(\mathsf H\), so
\eqref{eq:strong-neighborhood-from-hilbert-ball} implies
\(\rho^{(J+1)}\in U_*\), a contradiction.  The tail is therefore trapped,
and letting \(J\to+\infty\) in \eqref{eq:ls-finite-tail-bound} proves
finite length.  The sequence is Cauchy in \(\mathsf H\) and converges to
\(\rho_*\), because the selected cluster subsequence converges to that same
point in \(\mathsf H\).
\end{proof}

\subsection{Verification for entropic GW}

For compact GW, the required infinite-dimensional gradient inequality is not
an extra numerical assumption.  It follows from analyticity of entropy and
compactness of the integral operator.

\[
 \mathsf T_\infty
 =
 \mathsf T\cap L^\infty(\mathfrak m),
 \qquad
 \norm{h}_{\mathsf T_\infty}=\norm{h}_{L^\infty(\mathfrak m)}.
\]
This is a closed Banach subspace because the two marginal operators are
bounded on \(L^\infty\).

\begin{lemma}[The constrained entropy Hessian is invertible]
\label{lem:entropy-hessian-isomorphism}
Let \(\rho\in\mathcal A\cap L^\infty(\mathfrak m)\) satisfy
\(0<c\leq\rho\leq C<+\infty\).  Then
\begin{equation}\label{eq:constrained-entropy-hessian}
 \mathsf A_\rho:\mathsf T_\infty\longrightarrow\mathsf T_\infty,
 \qquad
 \mathsf A_\rho h
 =
 \mathsf P_{\mathsf T}\left(\frac h\rho\right),
\end{equation}
is a bounded isomorphism.
\end{lemma}

\begin{proof}
Boundedness is immediate.  If \(\mathsf A_\rho h=0\), then
\eqref{eq:coupling-tangent-projection} shows that
\(h/\rho=a(x)+b(y)\) for some bounded \(a,b\).  Since \(h\) has zero
marginals,
\[
 \int\frac{h^2}{\rho}\dd\mathfrak m
 =
 \int h(x,y)\bigl(a(x)+b(y)\bigr)\dd\mathfrak m(x,y)
 =0,
\]
and hence \(h=0\).  Thus \(\mathsf A_\rho\) is injective.

To prove surjectivity, fix \(g\in\mathsf T_\infty\).  Introduce the two
Markov operators
\[
 (Rb)(x)=\int\rho(x,y)b(y)\dd\beta(y),
 \qquad
 (R^*a)(y)=\int\rho(x,y)a(x)\dd\alpha(x),
\]
and the bounded functions
\[
 r_{\mathcal X}(x)=\int\rho(x,y)g(x,y)\dd\beta(y),
 \qquad
 r_{\mathcal Y}(y)=\int\rho(x,y)g(x,y)\dd\alpha(x).
\]
We seek \(h=\rho(g+a+b)\).  Its two marginals vanish precisely when
\begin{equation}\label{eq:entropy-hessian-potential-system}
 a+Rb=-r_{\mathcal X},
 \qquad
 R^*a+b=-r_{\mathcal Y}.
\end{equation}
Eliminating \(b\) gives
\begin{equation}\label{eq:entropy-hessian-schur-complement}
 (I-RR^*)a=-r_{\mathcal X}+Rr_{\mathcal Y}.
\end{equation}
Denote the right-hand side by \(r\).  It has zero \(\alpha\)-mean.  The
Markov operator \(P=RR^*\) has transition density
\[
 q(x,x')
 =
 \int\rho(x,y)\rho(x',y)\dd\beta(y)
 \geq c^2
\]
with respect to \(\alpha\).  Its rows integrate to one and, because
\(q\geq c^2\), it admits the minorization
\(q(x,\cdot)\alpha\geq c^2\alpha\).  Hence, for every bounded \(a\),
\begin{equation}\label{eq:entropy-hessian-dobrushin-contraction}
 \operatorname{osc}(Pa)
 \leq
 (1-c^2)\operatorname{osc}(a),
 \qquad
 \operatorname{osc}(a)=\operatorname*{ess\,sup}a-
                         \operatorname*{ess\,inf}a.
\end{equation}
Indeed, subtracting the common component \(c^2\alpha\) leaves a Markov
kernel of mass \(1-c^2\).  Thus \(P\) is a strict contraction on
\(L^\infty(\alpha)/\R\one\), and the Neumann series
\(\sum_{n\geq0}P^nr\) converges in that quotient.  Since \(P\) preserves
\(\alpha\) and \(r\) has zero mean, its zero-mean representative is a
bounded solution \(a\) of
\eqref{eq:entropy-hessian-schur-complement}.  Moreover,
\(\norm{r_{\mathcal X}}_\infty,\norm{r_{\mathcal Y}}_\infty
\leq\norm g_\infty\), so \eqref{eq:entropy-hessian-dobrushin-contraction}
bounds \(\norm a_\infty\) by a constant depending only on \(c\) times
\(\norm g_\infty\).

Set \(b=-r_{\mathcal Y}-R^*a\).  Then
\eqref{eq:entropy-hessian-potential-system} holds, so
\(h=\rho(g+a+b)\) lies in \(\mathsf T_\infty\).  Finally,
\(\mathsf P_{\mathsf T}(h/\rho)=g\), because the projection annihilates
\(a(x)+b(y)\), and
\(\norm h_\infty\leq C_\rho\norm g_\infty\).  This proves surjectivity and
boundedness of \(\mathsf A_\rho^{-1}\).
\end{proof}

\begin{proposition}[{\L}ojasiewicz--Simon inequality for entropic GW]
\label{prop:ls-entropic-gw}
Under the assumptions of
Proposition~\ref{prop:uniform-positive-gw-updates}, every strictly positive
stationary density
\(\rho_*\in\mathcal A\cap C(\mathcal X\times\mathcal Y)\) satisfies
\eqref{eq:lojasiewicz-simon-density} with \(\Phi=\Phi_\eps\),
\(G=G_\eps\), \(\mathsf B=\mathsf G=\mathsf T_\infty\).
\end{proposition}

\begin{proof}
Set \(\mathsf X=\mathsf T_\infty\).  The \(L^2(\mathfrak m)\) pairing defines
a continuous embedding
\[
 \jmath:\mathsf X\longrightarrow\mathsf X^*,
 \qquad
 \jmath(q)[h]=\int hq\dd\mathfrak m.
\]
It is definite because
\(\jmath(q)[q]=\norm q_{L^2(\mathfrak m)}^2>0\) for \(q\neq0\).
On the open set
\[
 U_*=
 \left\{
 h\in\mathsf X:
 \operatorname*{ess\,inf}(\rho_*+h)>0
 \right\},
\]
define
\[
 \mathcal E_*(h)=\Phi_\eps(\rho_*+h),
 \qquad
 \mathcal M_*(h)=G_\eps(\rho_*+h)\in\mathsf X.
\]
Stationarity and strict positivity give \(\mathcal M_*(0)=0\).  The
quadratic interaction is analytic on \(L^\infty\), and the Nemytskii map
\(\rho\mapsto\log\rho\) is real analytic on every uniformly positive
\(L^\infty\)-ball.  Hence \(\mathcal E_*\) and \(\mathcal M_*\) are real
analytic.  Moreover, for \(h,k\in\mathsf X\), self-adjointness of
\(\mathsf P_{\mathsf T}\) in \(L^2(\mathfrak m)\) and
\(k\in\mathsf T\) give
\[
 D\mathcal E_*(h)[k]
 =
 \int k\,\mathcal M_*(h)\dd\mathfrak m,
\]
so \(\mathcal M_*\) is a gradient map for \(\mathcal E_*\) with values in
\(\mathsf X\subset\mathsf X^*\).  Its derivative at the critical point
\(h=0\) is
\begin{equation}\label{eq:entropic-gw-linearized-gradient}
 D\mathcal M_*(0)h
 =
 2\mathsf P_{\mathsf T}Kh
 +
 \eps\mathsf A_{\rho_*}h,
 \qquad h\in\mathsf T_\infty.
\end{equation}

Because \(L\) is uniformly continuous on the compact product, the image by
\(K\) of the \(L^\infty\)-unit ball is uniformly bounded and
equicontinuous in \(C(\mathcal X\times\mathcal Y)\).  Arzel\`a--Ascoli and
the boundedness of \(\mathsf P_{\mathsf T}\) on continuous functions show
that \(\mathsf P_{\mathsf T}K:\mathsf X\to\mathsf X\) is compact.  By
Lemma~\ref{lem:entropy-hessian-isomorphism},
\eqref{eq:entropic-gw-linearized-gradient} is an isomorphism plus a compact
operator and is Fredholm of index zero.  The analytic-gradient-map
{\L}ojasiewicz--Simon theorem of \citet{FeehanMaridakis2020} applies with
the definite embeddings
\(\mathsf X=\widetilde{\mathsf X}=\mathsf T_\infty
\hookrightarrow\mathsf X^*\).  It yields constants
\(Z,\sigma>0\) and \(\vartheta\in[1/2,1)\) such that
\[
 \norm{G_\eps(\rho)}_\infty
 \geq
 Z\left|\Phi_\eps(\rho)-\Phi_\eps(\rho_*)\right|^\vartheta
 \quad
 \text{for }\norm{\rho-\rho_*}_\infty<\sigma,
\]
for feasible \(\rho\) in the positive chart.  Setting
\(\theta_*=1-\vartheta\in(0,1/2]\) and \(c_*=Z^{-1}\) gives
\eqref{eq:lojasiewicz-simon-density}.  This is the precise
infinite-dimensional replacement for the finite-dimensional KL property;
see also \citet{Chill2003}.
\end{proof}

\begin{theorem}[Full convergence of entropic alternating GW on measures]
\label{thm:entropic-gw-full-convergence}
Assume that \(\mathcal X,\mathcal Y\) are compact metric spaces, replace them
by the supports of \(\alpha,\beta\), and assume that the symmetric GW kernel
\(L\) is continuous and negative on feasible differences.  For every fixed
\(\eps>0\), initialize \eqref{eq:measure-alternating-update} with a positive
continuous coupling density, for instance
\(\rho^{(0)}=1\).  Then there is a continuous stationary fixed-point density
\(\rho_*\), satisfying \(c_\eps\leq\rho_*\leq C_\eps\), such that
\begin{equation}\label{eq:entropic-gw-full-convergence}
 \sum_{\ell=0}^{+\infty}
 \norm{\rho^{(\ell+1)}-\rho^{(\ell)}}_{L^2(\mathfrak m)}
 <+\infty,
 \qquad
 \rho^{(\ell)}\longrightarrow\rho_*
 \quad\text{uniformly}.
\end{equation}
Equivalently, \(\pi^{(\ell)}=\rho^{(\ell)}\mathfrak m\) converges to the
stationary coupling \(\pi_*=\rho_*\mathfrak m\).  The theorem does not claim
that \(\pi_*\) is a global GW minimizer, and its constants are not uniform as
\(\eps\downarrow0\).
\end{theorem}

\begin{proof}
Proposition~\ref{prop:exact-entropy-descent} and
\eqref{eq:kl-fisher-equivalence} give
\eqref{eq:abstract-sufficient-decrease} after the first, immaterial update;
explicitly one may take \(a=\eps/(2C_\eps)\).
Proposition~\ref{prop:block-residual-estimate} gives
\eqref{eq:abstract-relative-error}.  The iterates are relatively compact in
the uniform topology by Proposition~\ref{prop:uniform-positive-gw-updates},
\(\Phi_\eps\) is continuous on their uniformly positive closure, and every
cluster point is stationary.  Proposition~\ref{prop:ls-entropic-gw}
and Theorem~\ref{thm:measure-ls-convergence-principle} apply to the shifted
iterates \(h^{(\ell)}=\rho^{(\ell)}-1\), the shifted functional
\(\widehat\Phi_\eps(h)=\Phi_\eps(1+h)\), the residual
\(\widehat G_\eps(h)=G_\eps(1+h)\), and
\[
 \mathsf B=\mathsf T_\infty,
 \qquad
 \mathsf H=\mathsf T\subset L^2(\mathfrak m),
 \qquad
 \mathsf G=\mathsf T_\infty,
\]
They prove finite length and \(L^2\)-convergence to a stationary density
\(\rho_*\).  To upgrade the convergence, suppose that uniform convergence
failed.  A subsequence would remain a fixed positive uniform distance from
\(\rho_*\); relative compactness in
\(C(\mathcal X\times\mathcal Y)\) would give a uniformly convergent
subsequence.  Its limit must equal \(\rho_*\) in \(L^2(\mathfrak m)\), and
full support of \(\mathfrak m\) identifies the two continuous
representatives, a contradiction.
\end{proof}

\subsection{Extension beyond quadratic GW}

Only two algorithmic estimates were used: each exact block minimization must
produce a quadratic Bregman decrease, and changing the other block must
control the residual gradient.  This gives a genuinely biconvex
measure-space statement.

\begin{theorem}[Analytic biconvex alternating minimization]
\label{thm:analytic-biconvex-alternating}
Let \(\mathsf X=x_0+\mathsf X_0\) and
\(\mathsf Y=y_0+\mathsf Y_0\) be affine spaces modeled on Banach spaces
\(\mathsf X_0,\mathsf Y_0\), continuously and injectively embedded in
Hilbert spaces \(\mathsf H_{\mathsf X},\mathsf H_{\mathsf Y}\).  Let
\(\Bcal\) be analytic on a neighborhood of
\(\mathsf X\times\mathsf Y\) and convex in each block.  Consider the exact
updates
\begin{equation}\label{eq:general-two-block-updates}
 y^{(\ell+1)}
 \in\argmin_y\Bcal(x^{(\ell)},y),
 \qquad
 x^{(\ell+1)}
 \in\argmin_x\Bcal(x,y^{(\ell+1)}).
\end{equation}
All derivatives below are restricted to the model spaces
\(\mathsf X_0\) and \(\mathsf Y_0\), and all product spaces carry their
standard Euclidean product norms.
Assume that the iterates have compact closure in the product Banach topology,
that \(\Bcal\) is continuous there, and that the following properties hold
on an invariant neighborhood of this closure:
\begin{enumerate}[label=\textnormal{(\roman*)},leftmargin=2.2em]
\item both block minimizers are unique;
\item for some \(a>0\), each block update decreases \(\Bcal\) by at least
  \(a/2\) times the squared norm of that block increment in its Hilbert
  tangent space;
\item the constrained cross derivative is Lipschitz: for some \(b>0\),
  \[
   \norm{
    D_y\Bcal(x,y)-D_y\Bcal(\widetilde x,y)
   }_{\mathsf Y_0^*}
   \leq
   b\norm{x-\widetilde x}_{\mathsf H_{\mathsf X}};
  \]
\item at every stationary cluster point \(z_*\), there are
  \(c_*>0\), \(\theta_*\in(0,1/2]\), and a product-Banach neighborhood
  on which
  \[
   |\Bcal(z)-\Bcal(z_*)|^{1-\theta_*}
   \leq
   c_*\norm{D\Bcal(z)}_{\mathsf X_0^*\times\mathsf Y_0^*}.
  \]
\end{enumerate}
Then the full-cycle sequence
\(z^{(\ell)}=(x^{(\ell)},y^{(\ell)})\) has finite length in
\(\mathsf H_{\mathsf X}\times\mathsf H_{\mathsf Y}\) and converges to one
stationary point of \(\Bcal\).
\end{theorem}

\begin{proof}
Adding the two block-decrease inequalities gives
\begin{equation}\label{eq:abstract-biconvex-cycle-decrease}
 \Bcal(z^{(\ell)})-\Bcal(z^{(\ell+1)})
 \geq
 \frac a2
 \left(
  \norm{x^{(\ell+1)}-x^{(\ell)}}_{\mathsf H_{\mathsf X}}^2
  +
  \norm{y^{(\ell+1)}-y^{(\ell)}}_{\mathsf H_{\mathsf Y}}^2
 \right).
\end{equation}
Continuity on the compact closure bounds \(\Bcal\) below, so
\eqref{eq:abstract-biconvex-cycle-decrease} makes both block increments tend
to zero.

At the end of a cycle,
\(D_x\Bcal(x^{(\ell+1)},y^{(\ell+1)})=0\).  Optimality of the preceding
\(y\)-update gives \(D_y\Bcal(x^{(\ell)},y^{(\ell+1)})=0\).  Therefore
assumption~(iii) yields the relative-error estimate
\begin{align}
 \norm{D\Bcal(z^{(\ell+1)})}_{\mathsf X_0^*\times\mathsf Y_0^*}
 &=
 \norm{
  D_y\Bcal(x^{(\ell+1)},y^{(\ell+1)})
  -D_y\Bcal(x^{(\ell)},y^{(\ell+1)})
 }_{\mathsf Y_0^*}
 \notag\\
 &\leq
 b\norm{x^{(\ell+1)}-x^{(\ell)}}_{\mathsf H_{\mathsf X}}.
 \label{eq:abstract-biconvex-relative-error}
\end{align}
Thus the constrained gradient tends to zero.  Compactness and continuity of
\(D\Bcal\) make every cluster point stationary, so assumption~(iv) applies
at each of them.  Equations
\eqref{eq:abstract-biconvex-cycle-decrease} and
\eqref{eq:abstract-biconvex-relative-error} are precisely the
sufficient-decrease and relative-error estimates of
Theorem~\ref{thm:measure-ls-convergence-principle}, with
\(\mathsf B=\mathsf X_0\times\mathsf Y_0\),
\(\mathsf H=\mathsf H_{\mathsf X}\times\mathsf H_{\mathsf Y}\), and
\(\mathsf G=\mathsf X_0^*\times\mathsf Y_0^*\).  That theorem gives finite
length and convergence to one stationary point.
\end{proof}

For integral energies, assumption~(ii) is naturally furnished by a uniformly
coercive Legendre or entropy Hessian, while assumption~(iii) follows from a
bounded cross-kernel.  Analyticity and a Fredholm constrained Hessian imply
assumption~(iv) by the same argument as
Proposition~\ref{prop:ls-entropic-gw}.  The theorem therefore covers
bilinear and biquadratic interactions, and more general smooth biconvex
integral functionals, over spaces of measures.  It is the
infinite-dimensional counterpart of classical block-coordinate convergence
theory \citep{Tseng2001}.

\begin{remark}[Why positive regularization matters]
For \(\eps=0\), a GW block update is a linear problem and can have many
minimizers.  The Bregman term in
\eqref{eq:exact-entropic-gw-descent}, the positive density box, and the
coercive part of \eqref{eq:entropic-gw-linearized-gradient} all disappear.
Monotonicity from Proposition~\ref{prop:measure-descent} remains, but a
selection rule can cycle on a constant-value face.  Full convergence of the
unregularized iteration therefore requires a separate uniqueness, proximal,
or active-set hypothesis.
\end{remark}

\section{Practical conclusions}

The measure formulation separates properties that are often conflated.
\begin{center}
\renewcommand{\arraystretch}{1.18}
\begin{tabularx}{\linewidth}{@{}>{\bfseries\raggedright\arraybackslash}p{.31\linewidth}X@{}}
\toprule
Question & Correct statement \\
\midrule
What is duplicated? &
A measure \(\pi\in\Kcal\): the diagonal form \(\Ee(\pi,\pi)\) becomes
\(\Ee(\pi,\xi)\). \\
Why is the relaxation tight? &
Because \(\Ee\) is nonpositive on \(\Span(\Kcal-\Kcal)\), as exposed by
\eqref{eq:measure-polarization-identity}. \\
Why does this apply to GW? &
Differences of couplings have zero marginals, and negative-type geometry gives
the square identity \eqref{eq:measure-gw-negative-square}. \\
What is one block update? &
Linear OT for \(\eps=0\), entropic OT with strength \(\eps/2\) for
\(\eps>0\). \\
What converges on compact spaces? &
For entropic GW, the whole sequence converges uniformly to a stationary
positive coupling density and has finite Fisher-equivalent \(L^2\) length. \\
When does the whole sequence converge? &
Whenever sufficient Bregman decrease, a relative-gradient error bound,
precompactness, and a {\L}ojasiewicz--Simon inequality hold; these properties
are verified here for compact entropic GW. \\
Does it find the global GW minimum? &
Not in general.  The algorithm remains a local method for a nonconvex diagonal
objective. \\
\bottomrule
\end{tabularx}
\end{center}

\section{Conclusion}

The basic GW iteration is an instance of a general measure-level construction.
A negative symmetric bilinear form can be polarized without changing its
global diagonal minimum, while each variable of the polarized problem enters
linearly.  Squared GW inherits the required negativity from the Hilbert
embeddings of negative-type distances.  Entropy fits as a separable strictly
convex penalty and turns every frozen-measure problem into a Sinkhorn solve.
This common structure explains tightness, the factor \(\eps/2\), descent, and
stationarity in one argument.  More importantly, it exposes the right
convergence geometry.  The exact decrease is a relative entropy, its local
quadratic form is the Fisher metric, and the linearized constrained gradient
is an entropy isomorphism perturbed by a compact GW operator.  The resulting
{\L}ojasiewicz--Simon argument proves finite length and convergence directly
for measure-valued iterates, without passing through a matrix norm.  The same
proof architecture extends to analytic biconvex integral energies with
coercive block geometry and controlled cross derivatives.

\small
\setlength{\bibsep}{3pt plus 1pt minus 1pt}
\bibliographystyle{plainnat}
\bibliography{references}

\end{document}
